The results of our experiments provide strong support for the Discord OpenClaw framework as an effective approach to building safe, proactive AI agents in communication platforms. Several important insights emerge from our findings.

First, the substantial improvements in engagement metrics demonstrate that proactive behavior is valued by users when implemented appropriately. The fact that users engaged more frequently and for longer periods with proactive suggestions indicates that the system successfully anticipated user needs without being intrusive. This addresses a key limitation of purely reactive agents that miss opportunities to assist.

Second, the dramatic improvement in safety metrics with the approval workflow challenges the assumption that safety mechanisms necessarily reduce usability. Despite requiring an extra step for tool execution, users did not find the process burdensome, and the significant reduction in harmful actions suggests that users appreciate having control over potentially impactful operations. The minimal impact on task completion time indicates that the approval workflow introduces little friction for beneficial actions while effectively preventing harmful ones.

Third, the substantial reduction in hallucination rates through external knowledge integration highlights the importance of grounding language model responses in verified information. The improvement in citation quality additionally shows that users value being able to verify the sources of information provided by the agent. This addresses a fundamental limitation of language models that can compromise their usefulness in knowledge-intensive contexts.

Fourth, the ablation study reveals interesting insights about the relative importance of each component. While we expected proactive behavior to have the largest impact, the approval workflow component showed nearly equal importance, suggesting that users place high value on having control over the agent's actions. The external reasoning component, while showing smaller impact than the other two, still provided meaningful improvements in factual accuracy and hallucination reduction.

An unexpected finding was that user satisfaction did not show a significant increase with the approval workflow, despite the large improvements in safety and trust. This suggests that while users value safety and control, these factors may not directly translate to higher satisfaction scores in our measurement, or that the slight increase in task completion time offset potential satisfaction gains. Future work could explore different ways of measuring user satisfaction or alternative approval workflow designs.

The framework also demonstrates good scalability and adaptability. The same core principles could be applied to other communication platforms (Slack, Teams, etc.) or different external reasoning engines beyond NotebookLM. The modular design allows individual components to be adjusted or replaced based on specific use case requirements.

Limitations of our study include the relatively short duration of the engagement study (two weeks), which may not capture long-term effects or novelty decay. Additionally, our user sample, while diverse, may not represent all possible Discord community types. The experiments were also conducted in controlled environments, which may not fully reflect the complexity of real-world server dynamics.

Despite these limitations, our results provide strong evidence that combining proactive behavior, structured approval workflows, and external reasoning integration creates a synergistic effect that enhances both the usefulness and safety of AI agents in communication platforms. The framework offers a practical approach to deploying language model-based agents in real-world interactive environments where both helpfulness and safety are paramount.